\section*{Readings awaiting the scans}
\begingroup\small
These notes preserve unresolved alternatives without changing the retained reading. They are separate from the numbered apparatus of differences between the diplomatic and corrected texts. Each item needs the scan.

\medskip\noindent\textbf{SOL-EN-010 (folio 96).}\par
Retained: $f_{2*}(4)\circ(2)$. Proposed: $(f_{2*}(4)\circ(2))^{-1}$.
Both supplied French layers retain the un-inverted second path in the description following 3.5. The English follows that recorded reading. The added inverse is a new mathematical proposal; preserve both readings in the unresolved apparatus until the print and full diagram are rechecked.
\par\noindent\textit{Source: }\texttt{port/SOL\_NOTES\_5.md}. Status: needs the scan.

\medskip\noindent\textbf{SOL-EN-017 (folio 144).}\par
Retained: $(\{y\}\times Y,B)$. Proposed: $(\{x\}\times Y,B)$.
The supplied French source retains y in this pair. The English accurately carries it. Keep that recorded printed reading and expose the suggested x alongside it for scan review.
\par\noindent\textit{Source: }\texttt{port/SOL\_NOTES\_5.md}. Status: needs the scan.

\medskip\noindent\textbf{SOL-EN-026 (folio 160).}\par
Retained: the strict localization at $\bar z$. Proposed: the strict localization at $\bar s$.
Both French layers retain the reference to the localization at zbar in this sentence. Retain that reading with the suggested sbar in the unresolved notes for review against the scan.
\par\noindent\textit{Source: }\texttt{port/SOL\_NOTES\_5.md}. Status: needs the scan.

\medskip\noindent\textbf{SOL-EN-028 (folio 161).}\par
Retained: $R\Gamma(B'\cup D,i'_!(E')|B'-(B'\cup D))$. Proposed: $R\Gamma(B'-D,i'_!(E')|B'-D)$.
Both French layers preserve the printed union and nested set difference in this display. Keep the source reading, and record the proposed correction for scan review.
\par\noindent\textit{Source: }\texttt{port/SOL\_NOTES\_5.md}. Status: needs the scan.

\medskip\noindent\textbf{SOL-EN-029 (folio 170).}\par
Retained: $\underline{\operatorname{Hom}}_A^\bullet(E,E)\to A$. Proposed target: $A_\natural$.
French p.170 explicitly retains A in this sentence, while the nearby constructions use A\_natural. Preserve the recorded reading and show the proposed target in the unresolved apparatus.
\par\noindent\textit{Source: }\texttt{port/SOL\_NOTES\_5.md}. Status: needs the scan.

\medskip\noindent\textbf{SOL-EN-033 (folio 171).}\par
Retained reference: 5.3.6. Proposed reference: 5.3.5.
Both French layers cite 5.3.6 after (5.9.1). Keep the printed reference and record 5.3.5 as an unresolved alternative requiring the source page.
\par\noindent\textit{Source: }\texttt{port/SOL\_NOTES\_5.md}. Status: needs the scan.

\medskip\noindent\textbf{SOL-EN-046 (folio 185).}\par
Retained: $E_2=KA_{X_1}$. Proposed: $F_1=KA_{X_1}$.
The supplied corrected French also reads E\_2 here. The proposed substitution is preserved as an unresolved reading for scan review; the English keeps the French source.
\par\noindent\textit{Source: }\texttt{port/SOL\_NOTES\_5.md}. Status: needs the scan.

\medskip\noindent\textbf{SOL-EN-050 (folio 191).}\par
Retained: $D_c({}_AX)=D_c({}_{A[G]}X)$. Proposed left side: $D_c(G,{}_AX)$.
Both French layers preserve the left side without G. Retain it and record the proposed equivariant argument for scan review.
\par\noindent\textit{Source: }\texttt{port/SOL\_NOTES\_5.md}. Status: needs the scan.

\medskip\noindent\textbf{SOL-EN-052 (folio 195).}\par
Retained: $F^n\Lambda_X=0$ for $n\geq1$. Proposed bound: $n\geq2$.
Both French layers preserve the bound n>=1 alongside the separate F\textasciicircum{}1 definition. Retain the recorded printed bound and expose the suggested n>=2 for scan review.
\par\noindent\textit{Source: }\texttt{port/SOL\_NOTES\_5.md}. Status: needs the scan.

\medskip\noindent\textbf{SOL-EN-053 (folio 195).}\par
Retained target: $c_2^!(\Lambda_{X-U,X})$. Proposed target: $c_2^!(\Lambda_{U,X})$.
Both French layers retain the first target in the degree-one graded correspondence. Preserve that reading and show the alternative for scan review.
\par\noindent\textit{Source: }\texttt{port/SOL\_NOTES\_5.md}. Status: needs the scan.

\medskip\noindent\textbf{SOL-EN-056 (folio 197).}\par
Retained: every point $x$ of $G$. Proposed: every point $x$ of $C$.
The supplied French also reads "pour tout point x de G". Preserve it and show C as an alternative for review against the page.
\par\noindent\textit{Source: }\texttt{port/SOL\_NOTES\_5.md}. Status: needs the scan.

\medskip\noindent\textbf{SOL-EN-057 (folio 197).}\par
Retained: $f_*M$. Proposed: $f_*L$.
Both French layers retain f\_*M in Proposition 6.19. Keep the recorded reading and preserve the proposed L alongside it in the unresolved apparatus.
\par\noindent\textit{Source: }\texttt{port/SOL\_NOTES\_5.md}. Status: needs the scan.

\medskip\noindent\textbf{SOL-EN-062 (folio 206).}\par
Retained: $\underline{\operatorname{Hom}}_A(M,\operatorname{Hom}_{C}(X,Y))$. Proposed outer symbol: $\operatorname{Hom}_A$.
French p.206 has the underlined outer Hom in the supplied source, as does the ported English. Preserve this explicit source notation and record the proposed plain reading for scan review.
\par\noindent\textit{Source: }\texttt{port/SOL\_NOTES\_5.md}. Status: needs the scan.

\medskip\noindent\textbf{SOL-EN-073 (folio 371).}\par
Retained punctuation: ``cf. SGA 3 I 4.7.2). Proposed punctuation: ``cf. SGA 3 I 4.7.2.
Both supplied French layers contain the unmatched closing parenthesis at the same point. Retain the recorded punctuation and register the balanced alternative for inspection of the printed page.
\par\noindent\textit{Source: }\texttt{port/SOL\_NOTES\_5.md}. Status: needs the scan.

\medskip\noindent\textbf{SOL-EN-075 (folio 375).}\par
Retained: Weil formula 6.1 and local invariants in no.~5. Proposed: formula 5.2 and local invariants in no.~6.
Both French layers retain the two recorded references. Keep them in the translated body and preserve the suggested renumbering in the unresolved notes for scan review.
\par\noindent\textit{Source: }\texttt{port/SOL\_NOTES\_5.md}. Status: needs the scan.

\medskip\noindent\textbf{SOL-EN-076 (folio 379).}\par
Retained: defined in 4.1. Proposed: defined in 3.4.
French Proposition 3.6 explicitly says defini dans 4.1. Retain that source reference, with the suggested 3.4 recorded for inspection of the printed page.
\par\noindent\textit{Source: }\texttt{port/SOL\_NOTES\_5.md}. Status: needs the scan.

\medskip\noindent\textbf{SOL-EN-077 (folio 380).}\par
Retained reference: (4.3.1). Proposed reference: (3.7.1).
Both French layers give (4.3.1) in the sentence after (3.7.2). Preserve that reference and expose the alternative in the unresolved apparatus.
\par\noindent\textit{Source: }\texttt{port/SOL\_NOTES\_5.md}. Status: needs the scan.

\medskip\noindent\textbf{SOL-EN-086 (folio 475).}\par
Retained: an endomorphism of $K'$. Proposed: an endomorphism of $K'_\nu$.
The supplied French also says Kprime without the index nu in this sentence. Preserve that recorded reading, while retaining the indexed alternative for scan review.
\par\noindent\textit{Source: }\texttt{port/SOL\_NOTES\_5.md}. Status: needs the scan.

\medskip\noindent\textbf{REVIEW-5-01 (folio 299).}\par
Retained: ``faisceau normal (normal sheaf). Handwritten alternative: ``faisceau conormal (conormal sheaf).
The source records classify the thin added co as a reader annotation and preserve normal. N=J/J\textasciicircum{}2 and its dual occur explicitly in the same statement; replacing the noun is not required to recover its displayed formula. The printed-page image is outside the supplied range, so the handwritten insertion cannot be independently classified here.
\par\noindent\textit{Source: }\texttt{fixes/REVIEW\_5\_OPEN.jsonl}. Status: needs the scan.

\medskip\noindent\textbf{GLOBAL-85-LABELS (folio Exposé V).}\par
Retained: the source’s bold mixed-case statement labels. Recorded print alternative: DEFINITION 1.1.~- and the related capital labels.
The record gives example print labels but no complete per-label mapping, and the relevant scans are absent. Preserve the final source label forms; retain the stated printed alternatives as unresolved typography notes rather than infer one replacement across the exposé.
\par\noindent\textit{Source: }\texttt{conventions/GLOBAL\_ISSUES\_SGA.md}. Status: needs the scan.

\medskip\noindent\textbf{GLOBAL-85-ARROWS (folio 197-199).}\par
Retained: lower horizontal arrow labels below their arrows. Recorded print alternative: labels above those arrows.
The positional observation is volume-wide and was not itemized. The arrow sources, targets and labels are retained exactly; moving labels by a global rule would require the unavailable pages and a per-diagram check.
\par\noindent\textit{Source: }\texttt{conventions/GLOBAL\_ISSUES\_SGA.md}. Status: needs the scan.

\medskip\noindent\textbf{GLOBAL-89a (folio Exposé VII).}\par
Retained: bold mixed-case labels. Recorded print alternatives: capital labels with ``.-, a dash alone, or no stop or dash, according to the particular statement.
The source gives several different punctuation patterns, including one label without stop or dash. There is no complete per-label placement record, so no uniform replacement is justified from the attached text.
\par\noindent\textit{Source: }\texttt{conventions/GLOBAL\_ISSUES\_SGA.md}. Status: needs the scan.

\medskip\noindent\textbf{GLOBAL-89b (folio 281-295).}\par
Retained: the source’s itemized readings at each occurrence. Remaining alternatives for unitemized sites: plain $R$ or double-struck $\mathbb{R}$; the record names exceptions in $Rp_*$, $R_!p$, $Rf_*$ in the proof of 1.2.
The page checks applied occurrence-level font fixes, and the record explicitly identifies exceptions. Preserve those final occurrences and their macro scope; do not infer the unexamined ones from a single global definition.
\par\noindent\textit{Source: }\texttt{conventions/GLOBAL\_ISSUES\_SGA.md}. Status: needs the scan.

\medskip\noindent\textbf{GLOBAL-89c (folio 288-289).}\par
Retained: $\mathbb{P}$ and $\mathbb{E}$. Recorded print alternatives: typed $P$ and $E$ at 21 reported places.
The global note did not itemize the 21 sites. The unavailable pages are needed to map each claimed font reading to the exact occurrence. Preserve both stated readings in the unresolved notes.
\par\noindent\textit{Source: }\texttt{conventions/GLOBAL\_ISSUES\_SGA.md}. Status: needs the scan.

\medskip\noindent\textbf{GLOBAL-89e (folio Exposé VII).}\par
Retained: the capitalization checked at each occurrence. Pending alternatives at unitemized places: Module and module.
The source explicitly prohibits a blind pass and records capitalization page by page. Retain those final forms and list the general pending observation without changing unexamined prose.
\par\noindent\textit{Source: }\texttt{conventions/GLOBAL\_ISSUES\_SGA.md}. Status: needs the scan.

\medskip\noindent\textbf{GLOBAL-89f (folio Exposé VII, including 295).}\par
Retained: the source’s $\mathcal{O}$ occurrences. Recorded print alternatives: typed plain $O$ and hand-drawn script $\mathcal{O}_C$.
The supplied global note describes more than one glyph and gives no complete occurrence list. A new blanket replacement would discard that recorded variation. The final sources are retained and both glyph readings are exposed for later page inspection.
\par\noindent\textit{Source: }\texttt{conventions/GLOBAL\_ISSUES\_SGA.md}. Status: needs the scan.

\medskip\noindent\textbf{GLOBAL-VIII-CATEGORIES (folio 351-371).}\par
Retained: script $\mathcal{C},\mathcal{D},\mathcal{A},\mathcal{B}$. Recorded print alternatives: typed capitals $C,D,A,B$.
The unnumbered note identifies a glyph question without an occurrence-level list and points to the general convention. Preserve the final text; record the plain and script renderings for inspection rather than redefine four category macros globally.
\par\noindent\textit{Source: }\texttt{conventions/GLOBAL\_ISSUES\_SGA.md}. Status: needs the scan.

\medskip\noindent\textbf{GLOBAL-VIII-QUOTES (folio 351-371).}\par
Retained: each quotation mark in the final source. Recorded alternatives: straight typewriter quotes and guillemets; six pairs are mentioned.
The note says the two pairs on pp.354-355 were itemized. Preserve the final source quotation marks at those sites and elsewhere; the missing images prevent a new blanket six-pair transcription.
\par\noindent\textit{Source: }\texttt{conventions/GLOBAL\_ISSUES\_SGA.md}. Status: needs the scan.

\medskip\noindent\textbf{GLOBAL-VII-SET-DIFFERENCE (folio 301 and Exposé VII).}\par
Retained: the two itemized signs $X\mathbin{\dot{-}}Y$ and $X'\mathbin{\dot{-}}Y'$ on p.~301, and the existing signs elsewhere. Pending alternatives elsewhere: dotted minus and the existing ordinary set-difference sign.
The final source explicitly writes mathbin\{dot\{-\}\} in both p.301 complements. Other occurrences were not inspected in the supplied note, so retain their existing symbols and expose the pending alternative.
\par\noindent\textit{Source: }\texttt{conventions/GLOBAL\_ISSUES\_SGA.md}. Status: needs the scan.

\medskip\noindent\textbf{GLOBAL-98b (folio 441).}\par
Retained in both branches: $\underline{\operatorname{\mathbb{F}r}}_{Y/S}:Y\longrightarrow Y^{(p/S)}$. Unresolved alternative: $\underline{\operatorname{Fr}}_{Y/S}:Y\longrightarrow Y^{(p/S)}$.
GLOBAL-98 expressly leaves this occurrence open against the general glyph convention of item 78. MISPRINTS\_5 entry 572 records the underlined double-struck F but qualifies its glyph treatment. The p.441 scan is unavailable. Under the rule for unresolved readings outside the supplied scans, preserve that recorded printed glyph in all six readers and place the upright alternative in the nonvariant unresolved notes. The relative Frobenius morphism Y->Y\textasciicircum{}\{(p/S)\} and all other symbols of the display are unchanged.
\par\noindent\textit{Source: }\texttt{conventions/GLOBAL\_ISSUES\_SGA.md; fixes/MISPRINTS\_5.json, entry 572}. Status: needs the scan.
\endgroup
\clearpage
